\subsection{外乘法}

设 E 为希尔伯特空间。对每个整数 \(n \geqslant 0\)，以 \(\theta_{n}\) 表示从 \(\mathbf{T}^{n}(\mathbf{E})\) 到 \(\symbf{\Lambda}^{n}(\mathbf{E})\) 上的典范映射；于是

\[
\theta_ {n} (x _ {1} \otimes \ldots \otimes x _ {n}) = x _ {1} \wedge \ldots \wedge x _ {n}\tag{31}
\]

其中 \(x_{1}, \ldots, x_{n}\) 属于 E。设 \(p\) 与 \(q\) 为两个正整数；根据公式 (30) 与 (31)，我们有

\[
u \wedge v = \theta_ {p + q} \left(\frac {1}{(p !) ^ {1 / 2}} \psi_ {p} (u) \otimes \frac {1}{(q !) ^ {1 / 2}} \psi_ {q} (v)\right)\tag{32}
\]

其中 \(u \in \Lambda^p(\mathbf{E})\)，\(v \in \Lambda^q(\mathbf{E})\)。由于 \(\| \theta_n \| \leqslant (n!)^{1/2}\)，我们得到不等式

\[
\| u \wedge v \| \leqslant \left(\frac {(p + q) !}{p ! q !}\right) ^ {1 / 2} \| u \| \cdot \| v \|\tag{33}
\]

其中 \(u \in \symbf{\Lambda}^p(\mathrm{E})\)，\(v \in \symbf{\Lambda}^q(\mathrm{E})\)。因此，映射 \((u, v) \mapsto u \wedge v\) 可由连续性延拓为从 \(\hat{\symbf{\Lambda}}^p(\mathrm{E}) \times \hat{\symbf{\Lambda}}^q(\mathrm{E})\) 到 \(\hat{\symbf{\Lambda}}^{p + q}(\mathrm{E})\) 中的双线性映射，其范数至多等于 \(\left(\frac{(p + q)!}{p!q!}\right)^{1/2}\)（参看 V, p.~73, 练习 2）。我们仍把它记作 \((u, v) \mapsto u \wedge v\)。

命题 6. --- 设 E 为希尔伯特空间。我们有

\[
\| x \wedge u \| \leqslant \| x \| \cdot \| u \|\tag{34}
\]

其中 \(x\in \mathbf{E}\)，\(u\in \hat{\Lambda} (\mathrm{E})\)。

显然只需考虑 \(\| x\| = 1\) 的情形。

设 F 为 E 中由所有与 x 正交的向量组成的希尔伯特子空间。由于 E 是 F 与直线 K.x 的希尔伯特和，由 V, p.~34 的推论可知，映射 \((v,w)\mapsto v+x\wedge w\) 是从 \(\hat{\mathbf{A}}(\mathbf{F})\oplus\hat{\mathbf{A}}(\mathbf{F})\) 到 \(\hat{\mathbf{A}}(\mathbf{E})\) 上的希尔伯特空间同构。若 \(u=v+x\wedge w\)，其中 v, w 属于 \(\hat{\mathbf{A}}(\mathbf{F})\)，则有 \(x\wedge u=x\wedge v\)，从而 \(\|x\wedge u\|=\|v\|\leqslant(\|v\|^{2}+\|w\|^{2})^{1/2}=\|u\|\)。

推论 1. --- a) 设 \(x_1, \ldots, x_n\) 为希尔伯特空间 E 的元素。我们有

\[
\left\| x _ {1} \wedge \dots \wedge x _ {n} \right\| \leqslant \left\| x _ {1} \right\| \dots \left\| x _ {n} \right\|,\tag{35}
\]

其中等号仅当某个 \(x_{i}\) 为零或序列 \((x_{1},\dots,x_{n})\) 正交时才成立。

\begin{enumerate}
\def\labelenumi{\alph{enumi})}
\setcounter{enumi}{1}
\tightlist
\item
  设 \(x_{1}, \ldots, x_{n}, y_{1}, \ldots, y_{n}\) 为希尔伯特空间 E 的元素。我们有
\end{enumerate}

\[
\left| \det (\langle x _ {i}, y _ {j} \rangle) \right| \leqslant \| x _ {1} \| \dots \| x _ {n} \| \cdot \| y _ {1} \| \dots \| y _ {n} \|;\tag{36}
\]

若向量 \(x_{i}\) 与 \(y_{i}\) 都非零，则 (36) 中等号成立当且仅当 \((x_{1},\ldots ,x_{n})\) 与 \((y_{1},\dots,y_{n})\) 各自都是 \(\mathbf{E}\) 的同一向量子空间的正交基。

不等式 (35) 由命题 6 对 n 用归纳法得到；不等式 (36) 可由在 \(\hat{\Lambda}^{n}(E)\) 中应用 Cauchy-Schwarz 不等式以及 V, p.~33 的公式 (26) 推出。

假定序列 \((x_{1},\dots,x_{n})\) 正交。于是

\[
\left\| x _ {1} \wedge \dots \wedge x _ {n} \right\| ^ {2} = \det (\langle x _ {i} | x _ {j} \rangle) = \prod_ {i = 1} ^ {n} \left\| x _ {i} \right\| ^ {2}
\]

因为 \(\langle x_i|x_j\rangle = 0\)（对 \(i\neq j\)）。

现在假定向量 \(x_{1}, \ldots, x_{n}\) 都非零且不构成正交序列。由于 \(\| x_{1} \wedge \ldots \wedge x_{n} \|\) 对向量 \(x_{1}, \ldots, x_{n}\) 的依赖是对称的，我们可以假定 \(x_{1}\) 不与子空间 \(F\)（\(E\) 中由 \(x_{2}, \ldots, x_{n}\) 生成的）正交，且 \(F\) 不恰为 0。我们可以把 \(x_{1}\) 分解为 \(x_{1}' + y\) 的形式，其中 \(y \neq 0\) 属于 \(F\)，而 \(x_{1}'\) 与 \(F\) 正交；这时 \(\| x_{1}' \| < \| x_{1} \|\)。但是

\[
x _ {1} \wedge x _ {2} \wedge \ldots \wedge x _ {n} = x _ {1} ^ {\prime} \wedge x _ {2} \wedge \ldots \wedge x _ {n},
\]

因而

\[
\begin{array}{r l} \left\| x _ {1} \wedge \dots \wedge x _ {n} \right\| & \leqslant \left\| x _ {1} ^ {\prime} \right\| \left\| x _ {2} \right\| \dots \left\| x _ {n} \right\| \\ & <   \left\| x _ {1} \right\| \left\| x _ {2} \right\| \dots \left\| x _ {n} \right\|. \end{array}
\]

假定向量 \(x_{i}\) 与向量 \(y_{i}\) 都非零。关系式 (36) 中的等式等价于下列诸等式的合取

\[
\left| \left\langle x _ {1} \wedge \dots \wedge x _ {n} \mid y _ {1} \wedge \dots \wedge y _ {n} \right\rangle \right| = \| x _ {1} \wedge \dots \wedge x _ {n} \| \cdot \| y _ {1} \wedge \dots \wedge y _ {n} \|\tag{37}
\]

\[
\left\| x _ {1} \wedge \dots \wedge x _ {n} \right\| = \left\| x _ {1} \right\| \dots \left\| x _ {n} \right\|, \quad \left\| y _ {1} \wedge \dots \wedge y _ {n} \right\| = \left\| y _ {1} \right\| \dots \left\| y _ {n} \right\|.\tag{38}
\]

由证明的第一部分，等式 (38) 蕴含序列 \((x_{1},\ldots ,x_{n})\) 与 \((y_{1},\ldots ,y_{n})\) 各自都正交，这又蕴含 \(x_{1}\wedge \ldots \wedge x_{n}\neq 0\) 且 \(y_{1}\wedge \ldots \wedge y_{n}\neq 0\)。在这些条件下，关系式 (37) 蕴含存在一个标量 \(\lambda \neq 0\) 使 \(y_{1}\wedge \ldots \wedge y_{n} = \lambda x_{1}\wedge \ldots \wedge x_{n}\)（V, p.~3, 注 1）；换言之，\((x_{1},\dots,x_{n})\) 与 \((y_{1},\dots,y_{n})\) 是 E 的同一向量子空间的基（A, III, § 11, No.~13）。

推论 2. --- 设 \((a_{ij})_{1\leqslant i,j\leqslant n}\) 为一个复元素埃尔米特矩阵，其行列式为 D。假定不等式

\[
\sum_ {i, j = 1} ^ {n} a _ {i j} \overline {{z}} _ {i} z _ {j} \geqslant 0\tag{39}
\]

对所有复数 \(z_{1}, \ldots, z_{n}\) 成立。于是

\[
0 \leqslant \mathrm{D} \leqslant a _ {1 1} \dots a _ {n n}.\tag{40}
\]

假定 \(\mathbf{D}\) 非零；等式 \(\mathbf{D} = a_{11} \ldots a_{nn}\) 成立当且仅当 \(a_{ij} = 0\)（对所有 \(i \neq j\)）。

设 \(\Phi\) 为向量空间 \(\mathbf{C}^n\) 上由下式给出的埃尔米特型

\[
\Phi (\mathbf {z}, \mathbf {z} ^ {\prime}) = \sum_ {i, j = 1} ^ {n} a _ {i j} \overline {{{{z}}}} _ {i} z _ {j} ^ {\prime}
\]

其中 \(\mathbf{z} = (z_1, \dots, z_n)\) 与 \(\mathbf{z}' = (z_1', \dots, z_n')\) 属于 \(\mathbf{C}^n\)。根据假设，\(\Phi\) 是正的。

先假定 \(\Phi\) 是非退化的，即 D 非零。若 \((e_1, \ldots, e_n)\) 是 \(\mathbf{C}^n\) 的典范基，则 \(\Phi(\mathbf{e}_i, \mathbf{e}_j) = a_{ij}\)，于是推论 2 立即由推论 1 \(a)\)（令 \(x_i = \mathbf{e}_i\)）得出。

由于 \(a_{ii} = \Phi(e_i, e_i) \geqslant 0\)，当 \(D = 0\) 时不等式 (40) 也成立。

推论 3（« Hadamard 不等式 »）. --- 设 \((a_{ij})_{1 \leqslant i,j \leqslant n}\) 为一个复元素矩阵，其行列式为 D。令

\[
c _ {i} = \left(\sum_ {j = 1} ^ {n} | a _ {i j} | ^ {2}\right) ^ {1 / 2} \quad \text {for} \quad 1 \leqslant i \leqslant n,
\]

并令 \(m = \sup_{i,j}|a_{ij}|\)。于是我们有

\[
| \mathbf {D} | \leqslant c _ {1} \dots c _ {n} \leqslant m ^ {n}. n ^ {n / 2}.\tag{41}
\]

若 \(\mathbf{D} \neq 0\)，要使 \(|\mathbf{D}| = c_1 \ldots c_n\) 成立，必须且只需行 \(y_i = (a_{ij})_{1 \leqslant j \leqslant n}\)（矩阵 \((a_{ij})_{1 \leqslant i,j \leqslant n}\) 的行）是两两正交的向量。

给空间 \(C^{n}\) 赋以由下式定义的内积

\[
\langle \mathbf {z} | \mathbf {z} ^ {\prime} \rangle = \sum_ {i = 1} ^ {n} \overline {{z}} _ {i} z _ {i} ^ {\prime}.
\]

设 \((\mathbf{x}_{1},\ldots,\mathbf{x}_{n})\) 为 \(C^{n}\) 的典范基，\(y_{i}\) 为以 \(a_{ij}\)（\(1\leqslant j\leqslant n\)）为分量的向量。我们有 \(\|x_{i}\|=1\) 与 \(\|y_{i}\|=c_{i}\)（对 \(1\leqslant i\leqslant n\)）；此外 \(\langle x_{i}|y_{j}\rangle=a_{ji}\)。于是不等式 \(|D|\leqslant c_{1}\ldots c_{n}\) 及其取等条件都是 V, p.~35 推论 1 的特殊情形。显然有 \(c_{i}\leqslant m.n^{1/2}\)，因此 \(c_{1}\ldots c_{n}\leqslant m^{n}.n^{n/2}\)。

